\documentclass[a4paper]{article}

\usepackage{amsfonts}
\usepackage{amsmath}
\usepackage{array}
\usepackage{fullpage}
\usepackage[left=1cm,right=1cm,top=1cm,bottom=1cm,nohead,nofoot]{geometry}
\usepackage{graphicx}
\usepackage[latin1]{inputenc}

\newcommand{\vtopimage}[2]{
  \setlength\fboxsep{1pt}
  \setlength\fboxrule{0pt}
  \fbox{\includegraphics[width=#1]{#2}}
}
\newcommand{\partimage}[1]{\vtopimage{.1\textwidth}{#1}}
\newcommand{\partphoto}[1]{\partimage{photos/#1.jpg.eps}}
\newcommand{\partsymbol}[1]{\partimage{symbols/#1.svg.eps}}

\begin{document}

\setlength{\parindent}{0in}

\title{Chaosknoten - animated logo board 1.0 \\
       https://wiki.blinkenarea.org/index.php/Chaosknoten}

\author{Stefan Sch�rmans, BlinkenArea \\
        stefan@blinkenarea.org}

\maketitle

\section{Bauteile / Parts}
\label{sec:parts}

\begin{tabular}{rm{.1\textwidth}m{.1\textwidth}m{.3\textwidth}m{.3\textwidth}}
  \# & Photo & Symbol &
    Beschreibung & Description \\
  \hline
  \\
  42 & \partphoto{led} & \partsymbol{led} &
    LED: klein, teilweise transparent. Die Polung ist zu beachten.
    Die beiden gr�nen Punkte am rechten Rand kennzeichnen die Kathode.
    Im Symbol sind sie als zus�tzliche schwarze Punkte dargestellt. &
    LED: small, partly transparent. Pay attention to the polarity.
    The two green dots at the right side mark the cathode.
    In the symbol, they are shown as additional black dots. \\
  \\
  1 & \partphoto{resistor} & \partsymbol{resistor} &
    Widerstand: klein, Aufruck ``1002''. Die Polung ist unwichtig. &
    Resistor: small, label ``1002''. The polarity is not important. \\
  \\
  1 & \partphoto{capacitor} & \partsymbol{capacitor} &
    Kondensator: klein, braun, kein Aufruck. Die Polung ist unwichtig. &
    Capacitor: small, brown, no label. The polarity is not important. \\
  \\
  1 & \partphoto{transistor} & \partsymbol{transistor} &
    Transistor: kleines IC, 8 Pins. Die Polung ist zu beachten.
    Der kleine Punkt in der unteren linken Ecke markiert Pin 1.
    Im Symbol wird Pin 1 durch einen Halbkreis am linken Rand gekennzeichnet. &
    Transistor: small IC, 8 pins. Pay attention to the polarity.
    The small dot in the lower left corner marks pin 1.
    In the symbol, pin 1 is indicated by a half circle at the left. \\
  \\
  1 & \partphoto{microcontroller} & \partsymbol{microcontroller} &
    Mikrokontroller: gr��eres IC, 20 Pins. Die Polung ist zu beachten.
    Der kleine Punkt in der oberen linken Ecke markiert Pin 1.
    Im Symbol wird Pin 1 durch einen Halbkreis am oberen Rand gekennzeichnet. &
    Microcontroller: bigger IC, 20 pins. Pay attention to the polarity.
    The small dot in the upper left corner marks pin 1.
    In the symbol, pin 1 is indicated by a half circle at the top. \\
  \\
  1 & \partphoto{switch} & \partsymbol{switch} &
    Schalter: scharz, zwei kleine weisse Hebel, 4 Pins. Pin 1 ist dort, wo
    die Zahl 1 aufgedruckt ist. Im Symbol ist Pin 1 mit einem kleinen Kreis
    gekennzeichet. &
    Switch: black, two small white levers, 4 pins. Pin 1 is where the
    number 1 is printed onto the device. In the symbol, pin 1 is marked
    with a small circle. \\
  \\
  1 & \partphoto{battery_holder} & \partsymbol{battery_holder} &
    Batteriehalter: silbern, aus Blech gefertigt, Ausrichtung beachten.
    Auf das runde Pad sollte eine sehr d�nne Schicht L�tzinn
    aufgetragen werden. &
    Battery Holder: silver, made from sheet metal,
    pay attention to orientation.
    A very thin layer of solder should be added to the round pad. \\
\end{tabular}

\section{Platine / Board}

\begin{tabular}{m{.47\textwidth}m{.47\textwidth}}
  Die schematische Darstellung der Platine zeigt, an welche Position und in
  welcher Ausrichtung die Bauteile eingel�tet werden m�ssen. &
  The schemtic drawing of the board shows to what positions and in which
  orientation the parts have to be soldered. \\
  Es wird empfohlen, die Bauteile in der Reihenfolge ihrer Erw�hnung in
  Abschnitt~\ref{sec:parts} aufzul�ten.
  F�r eine kurze Einf�hrung in das SMD-L�ten, siehe
  Abschnitt~\ref{sec_smd_solder}. &
  It is recommended to solder the parts in the order they are mentioned
  in section~\ref{sec:parts}.
  For a brief introduction to SMD soldering, see
  section~\ref{sec_smd_solder}. \\
\end{tabular}

\subsection{Vorderseite / Front Side}

\begin{tabular}{cc}
  \vtopimage{.43\textwidth}{symbols/Chaosknoten_front.svg.eps}
  \vtopimage{.43\textwidth}{photos/Chaosknoten_front.jpg.eps}
\end{tabular}

\subsection{R�ckseite / Back Side}

\begin{tabular}{cc}
  \vtopimage{.43\textwidth}{symbols/Chaosknoten_back.svg.eps}
  \vtopimage{.43\textwidth}{photos/Chaosknoten_back.jpg.eps}
\end{tabular}

\subsection{Schalter / Switches}

\begin{tabular}{rm{.1\textwidth}m{.4\textwidth}m{.4\textwidth}}
  \# & ON/off & Beschreibung & Description \\
  \hline
  \\
  1 & ON  & eingeschaltet & turned on \\
    & off & ausgeschaltet & turned off \\
  2 & ON  & automatisch zur n�chsten Animation wechseln &
            switch to next animation automatically \\
    & off & aktuelle Animation weiter anzeigen &
            keep showing current animation \\
\end{tabular}

\section{Programmierung / Programming}

\begin{tabular}{m{.47\textwidth}m{.47\textwidth}}
  Der Mikrokontroller im Bausatz ist bereits programmiert.
  Daher kann die Platine nach Abschluss der L�tarbeiten direkt eingeschaltet
  werden. &
  The microcontroller in the kit is already programmed.
  Therefore, the board can be turned on directly after completing the
  sorldering.
  \\
  Wenn ein noch nicht programmierter Mikrokontroller eingesetzt wird,
  kann die Programmierung �ber die Debug-Kontakte in der linken unteren Ecke
  der R�ckseite durchgef�hrt werden.
  Dazu wird ein geeignetes Programmierger�t und ein Adapter ben�tigt.
  Die Pinbelegung des Adapters kann dem Schaltplan
  (siehe Abschnitt~\ref{sec:schematic}) entnommen werden. &
  If a non-programmed microncontroller is used, it can be programmed using
  the debug pads in the lower left corner of the back side.
  For this purpose, a suitable programmer and adapter are needed.
  The pinout of the adapter can be seen from the schematic
  (see section~\ref{sec:schematic}).
\end{tabular}

\section{Anhang / Appendix}

\subsection{Schaltplan / Schematic}
\label{sec:schematic}

\includegraphics[width=\textwidth]{symbols/schematic.svg.eps}

\subsection{St�ckliste / Part List}

\begin{tabular}{rlllll}
  \# & Beschreibung          & Description           & Bauform / Footprint & Wert / Value  & Symbol       \\
  \hline
   1 & Platine               & printed circuit board & 65mm x 50mm         & Chaosknoten   & -            \\
   1 & Batteriehalter        & battery holder        & SMD                 & -             & B1           \\
   1 & Batterie              & battery               & 20mm x 3.2mm        & CR2032        & f�r / for B1 \\
   1 & 2-pol. DIP-Schalter   & 2 pol DIP switch      & SMD                 & -             & S1           \\
   1 & Widerstand            & resistor              & SMD 0805            & 10kOhm        & R1           \\
   1 & Kondensator           & capacitor             & SMD 0805            & 100nF         & C1           \\
   1 & Transistor (N-FET)    & transistor (N-FET)    & SO-8                & IRF7402       & T1           \\
   1 & Mikrokontroller       & microcontroller       & SO-20               & ATtiny2313A   & IC1          \\
  42 & orange LED            & orange LED            & SMD                 & APL3015       & D1 - D42     \\
\end{tabular}

\subsection{SMD-L�ten / SMD Soldering}
\label{sec_smd_solder}

\begin{tabular}{rm{.1\textwidth}m{.37\textwidth}m{.37\textwidth}}
  1 & \partsymbol{smd_solder_start} &
    Das abgebildete SMD-Bauteil soll auf die Platine gel�tet werden
    (Ansicht von der Seite). &
    The pictured SMD part shall be soldered to the board
    (viewn from the side). \\
  \\
  2 & \partsymbol{smd_solder_pads} &
    Zun�chst wird das Bauteil noch einmal zur Seite gelegt. &
    First, the part is put aside. \\
  \\
  3 & \partsymbol{smd_solder_tin} &
    Mit L�tkolben und L�tzinn wird nun ein einziges Pad verzinnt.
    Falls es viele Pads gibt, w�hlt man ein Pad in einer Ecke.
    Die anderen Pads bleiben erst einmal frei. &
    Using soldering iron and solder, some solder is put onto a single pad.
    If there are many pads, a pad in a corner is selected.
    The other pads stay empty for now. \\
  \\
  3 & \partsymbol{smd_solder_part} &
    Mit dem L�tkolben wird nun das Zinn auf dem Pad wieder erhitzt.
    Gleichzeitig wird mit einer Pinzette in der anderen Hand
    das Bauteil an die vorgesehene Position gesetzt.
    Der L�tkolben wird von der L�tstelle entfernt.
    Wenn das L�tzinn erstarrt ist, wird das Bauteil losgelassen. &
    The solder on the pad is heated again with the soldering iron.
    At the same time, the part is placed to its designated position
    using some tweezers and the other hand.
    The soldering iron is removed from the solder joint.
    After the solder has solidified, the part is released. \\
  \\
  4 & \partsymbol{smd_solder_part} &
    Nun wird die Position des Bauteils noch einmal �berpr�ft.
    Insbesondere ist es wichtig, dass alle anderen Pins auf den vorgesehenen
    Pads liegen.
    Falls nicht, kann das Bauteil wie in Schritt 3 noch einmal neu
    positioniert werden, indem man die einzige L�tstelle noch einmal erhitzt. &
    Now the position of the part is verified.
    It is especially important that all other pins are exactly located at the
    designated pads.
    If this is not the case, the part can be re-positioned as in step 3 by
    heating the single solder joint again. \\
  \\
  5 & \partsymbol{smd_solder_done} &
    Wenn das Bauteil korrekt positioniert ist, werden alle anderen Pins mit
    Hilfe von L�tkolben und L�tzinn angel�tet.
    Falls das Bauteil viele Pins hat, l�tet man als zweites den
    Pin in der Ecke an, die dem zuerst angel�teten Pin gegen�ber liegt.
    Danach pr�ft man die Position des Bauteils noch einmal,
    bevor man die �brigen Pins anl�tet. &
    If the part is positioned correctly, all other pins are soldered using
    soldering iron and solder.
    In case the part has many pins, the pin in the corner opposite of the
    pin soldered first is soldered second.
    Afterwards, the position of the part is checked again,
    before soldering the remaining pins. \\
\end{tabular}

\section{Information}

\begin{tabular}{m{.47\textwidth}m{.47\textwidth}}
  Dieser Bausatz enth�lt Kleinteile und ist f�r kleine Kinder nicht geeignet. &
  This kit contains small parts and is not suitable for small children. \\
\end{tabular}

\center{
  Copyright (C) 2019 Stefan Sch�rmans, BlinkenArea

  Copyleft: GNU public license - https://www.gnu.org/copyleft/gpl.html
}

\end{document}

